
\documentclass[a4paper]{article}

% algumas packages para arrumar as tables co a margin:
% allows for temporary adjustment of side margins
\usepackage{changepage}

% provides filler text
\usepackage{lipsum}

% just makes the table prettier (see \toprule, \bottomrule, etc. commands below)
\usepackage{booktabs}

\usepackage[english, brazil]{babel} %Para traduzir os textos
\usepackage[utf8]{inputenc} %Para poder usar acentos
\usepackage[a4paper]{geometry} %Para ajustar a parte geometrica da folha
\geometry{verbose,tmargin=2.5cm,bmargin=2.5cm,lmargin=3cm,rmargin=3cm} % margem cute
\setlength{\parindent}{10mm}
\usepackage{wrapfig} %Biblioteca Matematica/Grafica
\usepackage{mathptmx} %Biblioteca Matematica/Grafica
\renewcommand{\ttdefault}{mathptmx} %Biblioteca Matematica/Grafica
\usepackage{amsmath} %Biblioteca Matematica/Grafica
\usepackage{amssymb} %Biblioteca Matematica/Grafica
\usepackage[11pt]{moresize}% different letters sizes
\usepackage{float}% enables accurate location of tables
\usepackage{caption}% to make personalized captions
\usepackage{graphicx} %Para inclusão de imagens
\usepackage{amsfonts}
\usepackage[T1]{fontenc}

%Packages do Lyx:
	\usepackage{units}

	%% Because html converters don't know tabularnewline
	\providecommand{\tabularnewline}{\\}
%Fim dos Packages do Lyx

\makeatletter

\title{Experimento 6:  Parte B \\ Calorimetria} % main title
\author{F 229 \\ \textsc{Grupo 1}}
\date{26 de Novembro, 2014}

\makeatother

\begin{document}
	% actually starts the document here
	\maketitle

	% members of the group


	\begin{center}
		\begin{tabular}{rr}
				\multicolumn{2}{c}{Integrantes}\tabularnewline
			                                  & \tabularnewline
			        Henrique Noronha Facioli  & RA: 157986 \tabularnewline
			        Guilherme Lucas da Silva  & RA: 155618 \tabularnewline
			          Beatriz Sechin Zazulla  & RA: 154779 \tabularnewline
			              Lucas Alves Racoci  & RA: 156331 \tabularnewline
			Isadora Sophia Garcia Rodopoulos  & RA: 158018 \tabularnewline
		\end{tabular}
	\par\end{center}

	%%Seria legal se alguém conseguir colocar os logos no footer da pagina, ou então no header...


	\begin{figure}[!ht]
		\vspace{8cm}
		\begin{minipage}[b]{0.45\linewidth}
			\includegraphics[scale=0.25]{logo-unicamp-name-line-blk-blk-0480.jpg} %
		\end{minipage}
		\hspace{0.5cm}
		\begin{minipage}[b]{0.45\linewidth}
			\includegraphics[scale=0.25]{logo-ifgw.png}
		\end{minipage}
	\end{figure}


	\newpage{}


	\section{Resumo}
		Neste experimento, estudamos o calorímetro, um instrumento usado na
		medição de processos que envolvam trocam de calor, ou seja, mudança
		de temperatura, estado ou qualquer outro processo que seja alterado
		pelo calor no sistema. Nesta segunda parte, estudamos e determinamos
		o calor específico de três metais distintos e o calor latente de fusão
		do gelo. Dadas equações já conhecidas do ramo da Calorimetria, obtivemos
		as equações finais que, no contexto do experimento, foram capazes
		de relacionar essas grandezas. Valendo-nos de tais leis e dos materiais
		disponíveis (calorímetro, termômetro e água a diferentes temperaturas),
		fomos capazes de encontrar os valores procurados, determinando diversas
		temperaturas de interesse (em estados iniciais e no equilíbrio térmico)
		e levando em consideração os erros inerentes a toda medição e ao uso
		de equipamentos reais.


	\section{Objetivo}
		Esta segunda parte do expermiento tem como objetivo o estudo e a determinação
		do calor específico de 3 metais ($c_{M1},c_{M2},c_{M3}$) e do calor
		latente de fusão do gelo ($L_{Gelo}$). Para tal, conta-se com um
		calorímetro, um termômetro e água a diferentes temperaturas de modo
		a encontrarmos as grandezas procuradas através de equações conhecidas,
		tais como:

		\begin{equation}
			Q=mc\Delta T
			\label{eq:calor transferido em funcao de c}
		\end{equation}


		\begin{equation}
			Q=mL
			\label{eq:calor latente de fusao}
		\end{equation}


		Onde, em \eqref{eq:calor transferido em funcao de c}, $Q$ corresponde
		ao calor transferido ou recebido por um material, dada uma massa $m$,
		uma variação de temperatura $\Delta T$ e seu calor específico, $c$.
		Em \eqref{eq:calor latente de fusao}, $Q$ também corresponde ao
		calor transferido por um material, porém no contexto de uma mudança
		de fase, em função de uma massa $m$ e um coeficiente de calor de
		latencia de variação do estado físico, que no nosso caso será a fusão,
		chamaremos esse coeficiente de $L$.


	\section{Procedimentos e coleta de dados}

		\subsection{Materiais utilizados}

			Para realizar o experimento necessitamos dos seguintes materiais:
			\begin{enumerate}
				\item Balança analítica 
				\item Calorímetro contendo:
					\begin{enumerate}
						\item Copo 
						\item Alça 
						\item Suporte 
						\item Isolamento 
						\item Tampa
					\end{enumerate}
				\item Três metais a serem identificados 
				\item Aquecedor 
				\item Gelo 
				\item Garrafa térmica 
				\item Béquer 
				\item Termômetro de mercúrio 
			\end{enumerate}

		\subsection{Procedimento}

			De posse dos materiais, primeiramente nós medimos a massa do copo
			do calorímetro vazio que será utilizada em todo o experimento para
			calcular a massa de água no calorímetro.

			Para a obtenção dessa massa de água, fizemos: $m_{a}=m_{c+\varnothing}-m_{c+a}$
			cujo erro fica: $\sigma_{m_{a}}=\sigma_{m}\cdot\sqrt{2}$, já que
			as massas $m_{c+\varnothing}$ e $m_{c+a}$ têm o mesmo erro instrumental,
			$\sigma_{m}$ .

			É importante ressaltar que, durante o experimento, o calor específico
			da água ($c_{a}$) foi considerado como 1 $\unitfrac{cal}{g\cdot^{\circ}C}$
			e o calor específico do gelo ($c_{g}$) como 0,5 $\unitfrac{cal}{g\cdot^{\circ}C}$


			\subsubsection{Parte 1: Obtenção do calor específico dos metais}

				Para cada um dos metais a serem identificados, medimos a massa do
				metal, a massa do copo do calorímetro com água, a temperatura dessa
				água no copo do calorímetro, que em alguns casos teve que ser resfriada
				com gelo, a temperatura do metal e a temperatura deequilíbrio.
	
				Para medir a temperatura do metal, primeiro colocamos água no béquer
				e esquentamos esta água até a fervura (a temperatura aqui não é relevante
				e portanto nem foi medida), à seguir, transferimos essa água do béquer
				para a garrafa térmica e colocamos o metal junto, esperamos até o
				que julgamos ser o equilíbrio entre o metal e esta água quente do
				béquer e medimos essa temperatura.
	
				Para medir a temperatura de equilíbrio, colocamos o metal já aquecido
				dentro do copo do calorímetro com água fria e colocamos esse conjunto
				no calorímetro. Esperamos novamente até o que julgamos ser o equilíbrio
				do sistema e medimos essa temperatura. Os valores obtidos para esses
				dados podem ser visualizados na tabela a seguir:
	
				\begin{table}[!ht]
					\caption{Valores medidos experimentalmente conforme expresso acima para a obtenção
					do calor especifico dos metais}
	
	
					\hspace{-1,7cm}%
					\begin{tabular}{|c|c|rl|rl|rl|}
						\hline 
						Dado Medido  & %
						\begin{tabular}{c}
					    	Símbolo \tabularnewline
						    Associado\tabularnewline
						\end{tabular} & 
						\multicolumn{2}{c|}{Para o metal $a$} & 
						\multicolumn{2}{c|}{Para o metal $b$} & 
						\multicolumn{2}{c|}{Para o metal $c$}\tabularnewline
						\hline 
						Massa do copo vazio  & $m_{c+\varnothing}$  & $\left(5,97\pm0,01\right)$  & $\cdot10^{1}\unit{g}$  & $\left(5,97\pm0,01\right)$  & $\cdot10^{1}\unit{g}$  & $\left(5,97\pm0,01\right)$  & $\cdot10^{1}\unit{g}$\tabularnewline
						\hline 
						Massa do copo com água  & $m_{c+a}$  & $\left(2,442\pm0,001\right)$  & $\cdot10^{2}\unit{g}$  & $\left(2,132\pm0,001\right)$  & $\cdot10^{2}\unit{g}$  & $\left(2,126\pm0,001\right)$  & $\cdot10^{2}\unit{g}$\tabularnewline
						\hline 
						Massa de água  & $m_{a}$  & $\left(1,825\pm0,001\right)$  & $\cdot10^{2}\unit{g}$  & $\left(1,535\pm0,001\right)$  & $\cdot10^{2}\unit{g}$  & $\left(1,529\pm0,001\right)$  & $\cdot10^{2}\unit{g}$\tabularnewline
						\hline 
						Temperatura da água  & $T_{a}$  & $\left(2,0\pm0,1\right)$  & $\cdot10^{1}\unit{^{\circ}C}$  & $\left(1,3\pm0,1\right)$  & $\cdot10^{1}\unit{^{\circ}C}$  & $\left(2,1\pm0,1\right)$  & $\cdot10^{1}\unit{^{\circ}C}$\tabularnewline
						\hline 
						Temperatura do metal  & $T_{q}$  & $\left(8,0\pm0,1\right)$  & $\cdot10^{1}\unit{^{\circ}C}$  & $\left(8,0\pm0,1\right)$  & $\cdot10^{1}\unit{^{\circ}C}$  & $\left(8,2\pm0,1\right)$  & $\cdot10^{1}\unit{^{\circ}C}$\tabularnewline
						\hline 
						Massa do metal  & $m_{M}$  & $\left(9,09\pm0,01\right)$  & $\cdot10^{1}\unit{g}$  & $\left(4,20\pm0,01\right)$  & $\cdot10^{1}\unit{g}$  & $\left(1,012\pm0,001\right)$  & $\cdot10^{2}\unit{g}$\tabularnewline
						\hline 
						Temperatura de equilíbrio  & $T_{e}$  & $\left(2,3\pm0,1\right)$  & $\cdot10^{1}\unit{\unit{^{\circ}C}}$  & $\left(2,0\pm0,1\right)$  & $\cdot10^{1}\unit{\unit{^{\circ}C}}$  & $\left(2,3\pm0,1\right)$  & $\cdot10^{1}\unit{\unit{^{\circ}C}}$\tabularnewline
						\hline 
					\end{tabular}
				\end{table}



			\subsubsection{Parte 2: Calor latente de fusão do gelo}

				Para realizar o cálculo que aproximo o calor latente de fusão do gelo,
				foi necessário, inicialmente, medir a massa da água e sua temperatura
				inicial e a massa de gelo e sua temperatura inicial.

				Em seguida, foi elaborado um sistema no calorímetro com a água e o
				gelo, que se finalizou com o gelo derretendo completamente - sendo
				medido, nesse instante, a sua temperatura de equilíbrio.

				Dipostos desses procedimentos, concluiu-se os dados suficientes para
				aproximar o calor latente de fusão do gelo. Os valores obtidos para
				esses dados podem ser visualizados na tabela a seguir:

				\newpage{}

						\begin{table}[!ht]
                			\caption{Valores medidos experimentalmente conforme expresso acima para a obtenção
                			do calor latente de fusão do gelo}
                
                
                			\centering{}%
                			\begin{tabular}{|c|c|rl|}
                				\hline 
                				Dado Medido  & %
                				\begin{tabular}{c}
                				Símbolo \tabularnewline
                				Associado\tabularnewline
                				\end{tabular} & \multicolumn{2}{c|}{Resultados}\tabularnewline
                				\hline 
                				Massa do copo vazio  & $m_{c+\varnothing}$  & $\left(5,97\pm0,01\right)$  & $\cdot10^{1}\unit{g}$ \tabularnewline
                				\hline 
                				Massa do copo com água  & $m_{c+a}$  & $\left(1,908\pm0,001\right)$  & $\cdot10^{2}\unit{g}$ \tabularnewline
                				\hline 
                				Massa de água  & $m_{a}$  & $\left(1,311\pm0,001\right)$  & $\cdot10^{2}\unit{g}$ \tabularnewline
                				\hline 
                				Temperatura da água  & $T_{a}$  & $\left(3,9\pm0,1\right)$  & $\cdot10^{1}\unit{^{\circ}C}$ \tabularnewline
                				\hline 
                				Massa do béquer vazio  & $m_{b+\varnothing}$  & $\left(5,82\pm0,01\right)$  & $\cdot10^{1}\unit{g}$ \tabularnewline
                				\hline 
                				Massa do béquer com gelo  & $m_{b+a}$  & $\left(8,57\pm0,01\right)$  & $\cdot10^{1}\unit{g}$ \tabularnewline
                				\hline 
                				Massa de gelo  & $m_{g}$  & $\left(2,75\pm0,01\right)$  & $\cdot10^{1}\unit{g}$ \tabularnewline
                				\hline 
                				Temperatura do gelo  & $T_{g}$  & $\left(0\pm1\right)$  & $\unit{^{\circ}C}$ \tabularnewline
                				\hline 
                				Temperatura de equilíbrio  & $T_{e}$  & $\left(2,30\pm0,05\right)$  & $\cdot10^{1}\unit{^{\circ}C}$ \tabularnewline
                				\hline 
                			\end{tabular}
                		\end{table}

	\section{Análise dos resultados}


		\subsection{Cálculo do calor específico dos metais}


			\subsubsection{Encontrando o calor específico de cada metal}

				Evoluindo da fórmula: $\sum Q=0$ - já que o sistema é quase fechado,
				têm-se, para cada metal:

				\begin{equation}
					Q_{calor\acute{\imath}metro}+Q_{\acute{a}gua}+Q_{metal}=0
				\end{equation}


				\[
					C\cdot\left(T_{e}-T_{a}\right)+c_{a}m_{a}\left(T_{e}-T_{a}\right)+c_{M}m_{M}\left(T_{e}-T_{q}\right)=0
				\]


				\[
					\left(T_{e}-T_{a}\right)\cdot\left(C+c_{a}m_{a}\right)=c_{M}m_{M}\left(T_{q}-T_{e}\right)
				\]


				\begin{equation}
					c_{M}=\cfrac{C+c_{a}m_{a}}{m_{M}}\cdot\cfrac{T_{e}-T_{a}}{T_{q}-T_{e}}\label{eq:c_m}
				\end{equation}


				Assim, fazendo $\Delta T_{e\rightarrow q}=T_{q}-T_{e}$ e $\Delta T_{a\rightarrow e}=T_{e}-T_{a}$
				o erro da equação \eqref{eq:c_m} é dado por:

				\begin{align*}
					\sigma_{c_{M}}^{2} = &
						\left(\sigma_{C}\cdot\cfrac{\partial c_{M}}{\partial C}\right)^{2}+
						\left(\sigma_{m_{a}}\cdot\cfrac{\partial c_{M}}{\partial m_{a}}\right)^{2}+
						\left(\sigma_{m_{M}}\cdot\cfrac{\partial c_{M}}{\partial m_{M}}\right)^{2}+
						\\ &
						\left(\sigma_{\Delta T_{e\rightarrow M}}\cdot\cfrac{\partial c_{M}}{\partial
							\left(\Delta T_{e\rightarrow q}\right)}\right)^{2}+
						\left(\sigma_{\Delta T_{a\rightarrow e}}\cdot\cfrac{\partial c_{M}}{\partial
							\left(\Delta T_{a\rightarrow e}\right)}\right)^{2}
                	\\
				    \sigma_{c_{M}}^{2} = & 
				        \left(\cfrac{T_{e}-T_{a}}{T_{q}-T_{e}}\cdot\cfrac{\sigma_{C}}{m_{M}}\right)^{2}+
						\left(\cfrac{T_{e}-T_{a}}{T_{q}-T_{e}}\cdot\cfrac{\sigma_{m_{a}}\cdot c_{a}}{m_{M}}\right)^{2}+
						\left(c_{M}\cdot\cfrac{\sigma_{m_{M}}}{m_{M}}\right)^{2}+ \\ &
						\left(c_{M}\cdot\cfrac{\sigma_{\Delta T_{e\rightarrow M}}}{\Delta T_{e\rightarrow q}}\right)^{2}+
						\left(c_{M}\cdot\cfrac{\sigma_{\Delta T_{a\rightarrow e}}}{\Delta T_{a\rightarrow e}}\right)^{2}
				\end{align*}
                
                Sendo $\sigma_{T}$ o erro instrumental, tem-se que 
                $\sigma_{\Delta T} = \sqrt{2\cdot \sigma_{T}^{2}}=\sigma_{T}\cdot \sqrt{2} $, 
                ou seja, $\sigma_{\Delta T_{e\rightarrow M}}= \sigma_{\Delta T_{a\rightarrow e}} = \sigma_{\Delta T}$, 
                assim, podemos fazer:

				\begin{equation}
					\sigma_{c_{M}} = 
					  \sqrt{
						\left(\cfrac{T_{e}-T_{a}}{T_{q}-T_{e}}\right)^{2}\cdot
						\cfrac{\sigma_{C}^{2}+\left(\sigma_{m_{a}}c_{a}\right)^{2}}{m_{M}^{2}}+
						c_{M}^{2}\cdot\left(
						\cfrac{\sigma_{m_{M}}}{m_{M}}+
						\cfrac{\sigma_{\Delta T}}{T_{q}-T_{e}}+
						\cfrac{\sigma_{\Delta T}}{T_{e}-T_{a}}\right)^{2}
					  }
				\end{equation}


				Lembrando da parte A desse experimento, têm-se que:

				\begin{equation}
					C=\left(1,8\pm0,8\right)\unit{\cdot10\cdot\frac{cal}{K}=\left(1,8\pm0,8\right)\cdot10\cdot}\frac{cal}{^{\circ}C}
				\end{equation}


				Assim obtém-se:

				\begin{table}[!ht]
					\caption{Valores calculados de $c_{M}\pm\sigma_{c_{M}}$ para cada metal}


					\centering{}%
					\begin{tabular}{|c|rl|}
						\cline{2-3} 
						\multicolumn{1}{c|}{} & \multicolumn{2}{c|}{$c_{M}\pm\sigma_{c_{M}}$}\tabularnewline
						\hline 
						Para o metal $a$  & $\left(1,2\pm0,4\right)$  & 
						$\cdot10^{-1}\unitfrac{cal}{g\cdot^{\circ}C}$\tabularnewline
						\hline 
						Para o metal $b$  & $\left(4,8\pm0,8\right)$  & 
						$\cdot10^{-1}\unitfrac{cal}{g\cdot^{\circ}C}$\tabularnewline
						\hline 
						Para o metal $c$  & $\left(6\pm3\right)$  & 
						$\cdot10^{-2}\unitfrac{cal}{g\cdot^{\circ}C}$\tabularnewline
						\hline 
					\end{tabular}
				\end{table}

			\subsubsection{Análise crítica dos resultados}

				Não nos foi fornecido o nome dos metais, por isso não é possível estabelecer
				com certeza se o experimento foi bem sucedido. Porém, o valor encontrado
				para o metal $a$ bate com o valor teórico esperado para 
				\begin{itemize}
					\item o ferro ($1,1\cdot10^{-1}\unitfrac{cal}{g\cdot^{\circ}C}$) e 
					\item o cobre($9,4\cdot10^{-2}\unitfrac{cal}{g\cdot^{\circ}C}$). 
				\end{itemize}
				O valor obtido para o metal $c$ também bate com o valor teórico esperado
				para a Prata ($5,6\cdot10^{-2}\unitfrac{cal}{g\cdot^{\circ}C}$).

				Porém não foi possível associar o valor encontrado para o metal $b$
				com o valor teórico esperado para nehum metal que se parecesse visualmente
				com ele.

				Assim, a partir das análises, foi observado que o valor do segundo metal não estava de acordo com o esperado. Isso pode ser explicado por: 

				\begin{enumerate}
					\item Quanto aos materiais: 
						\begin{enumerate}
							\item A imperfeição do calorímetro para trocar calor pode não ter sido capturada
							totalmente pelo erro na capacidade térmica; 
							\item Provavelmente o calorímetro usado na parte B do experimento não é
							exatamente o mesmo do usado na parte A, mas ainda que seja o mesmo
							calorímetro, o simples manuseio de outros grupos pode ter mudado sua
							capacidade térmica; 
							\item Um dos motivos para a suspeita de que o calorímetro não seja o mesmo,
							é porque notamos maior dificuldade para fechar este calo\'{r}imetro
							em comparação com o usado na parte A. 
						\end{enumerate}
					\item Quanto ao erro humano: 
						\begin{enumerate}
							\item É possível que o metal não tenha ficado tempo suficiente na água quente
							para que o equilíbrio fosse estabelecido e, portanto, a temperatura
							usada nos cálculos para o metal quente pode ter sido superestimada; 
							\item Em alguns casos a temperatura da água fria foi medida muito antes
							de o metal ser introduzido no sistema, logo esta água pode ter uma
							temperatura real no início do processo superior à medida. Assim, a
							temperatura usada nos cálculos talvez também esteja subestimada. 
						\end{enumerate}
					\item O processo de transferência do metal da garrafa térmica para o calorímetro
					pode tê-lo resfriado e, portanto, a temperatura medida da água quente
					pode ser relevantemente maior que real do metal, o que também comprometeria
					os cálculos; 
					\item A realização de uma única medida para a capacidade térmica, leva em
					conta apenas a propagação dos erros instrumentais do termometro e
					da balança, mas não leva em conta o erro estatístico que provavelmente
					se manifestaria em decorrência de erros aleatórios como a dissipação
					de calor por: 
						\begin{enumerate}
							\item Irradiação, já que o calorímetro não têm uma superfície suficientemente
							refletora; 
							\item Condução, já que o calorímetro não tem isolamento perfeito. 
						\end{enumerate}
					\item Também não foi considerado o erro associado ao calor específico da
					água. 
					\item Ao esquentar a água quente usada para calcular a capacidade térmica
					parte dela evapora, o que diminui a massa, portanto a medida obtida
					para a massa também fica superestimada. 
				\end{enumerate}

		\subsection{Cálculo do calor latente de fusão do gelo}

			\subsubsection{Encontrando o calor latente de fusão do gelo}

				Sabendo que: 
				\begin{equation}
					Q_{Calor\acute{i}metro}+
					Q_{\acute{A}gua}+Q_{Gelo}+
					Q_{Fus\tilde{a}o}+Q_{Gelo\rightarrow\acute{A}gua}=0 
					\Leftrightarrow
				\end{equation}


				\[
					C\cdot\left(T_{a}-T_{e}\right)+
					m_{a}\cdot c_{a}\cdot\left(T_{a}-T_{e}\right)+
					m_{g}c_{g}\left(T_{f}-T_{g}\right)+m_{g}L_{f}+
					m_{g}c_{a}\left(T_{e}-T_{a}\right)=0
					\Leftrightarrow
				\]


				\[
					L_{f} = - \cfrac{
										C\cdot\left(T_{e}-T_{a}\right)+
										m_{a}\cdot c_{a}\cdot\left(T_{e}-T_{a}\right)+
										m_{g}c_{g}\left(T_{f}-T_{g}\right)+
										m_{g}c_{a}\left(T_{e}-T_{f}\right)
									}
									{m_{g}}
				\]


				\begin{equation}
					L_{f} = \cfrac{\left(T_{a}-T_{e}\right)\cdot\left(C+m_{a}\cdot c_{a}\right)}{m_{g}}-c_{g}\cdot\left(T_{f}-T_{g}\right)-c_{a}\cdot\left(T_{e}-T_{f}\right)
				\end{equation}


				com erro dado por erro dado por:
				
	            \begin{align*}
	    			\sigma_{L_{f}}^{2} = &
	    				\left(\sigma_{T_{g}}\cdot\cfrac{\partial L_{f}}{\partial T_{g}}\right)^{2}+
	    				\left(\sigma_{T_{e}}\cdot\cfrac{\partial L_{f}}{\partial T_{e}}\right)^{2}+
	    				\left(\sigma_{T_{a}}\cdot\cfrac{\partial L_{f}}{\partial T_{a}}\right)^{2}+
	    				\left(\sigma_{T_{f}}\cdot\cfrac{\partial L_{f}}{\partial T_{f}}\right)^{2}+ \\ &
	    				\left(\sigma_{m_{a}}\cdot\cfrac{\partial L_{f}}{\partial m_{a}}\right)^{2}+
	    				\left(\sigma_{m_{g}}\cdot\cfrac{\partial L_{f}}{\partial m_{g}}\right)^{2}+
	    				\left(\sigma_{C}\cdot\cfrac{\partial L_{f}}{\partial C}\right)^{2} 
	    			\\
	    			= &
	    				\left(\sigma_{T_{g}}\cdot c_{g}\right)^{2}+
	    				\left(\sigma_{T_{e}}\cdot\cfrac{C+c_{a}\cdot\left(m_{a}+m_{g}\right)}{m_{g}}\right)^{2}+
	    				\left(\sigma_{T_{a}}\cdot\cfrac{C+m_{a}c_{a}}{m_{g}}\right)^{2}+
	    				\\ &
	    				\left(\sigma_{T_{f}}\cdot\left(c_{a}-c_{g}\right)\right)^{2}+
	    				\left(\sigma_{m_{a}}\cdot c_{a}\cdot\cfrac{T_{a}-T_{e}}{m_{g}}\right)^{2}+
	    				\left(\sigma_{m_{g}}\cdot\cfrac{\left(T_{a}-T_{e}\right)\cdot\left(C+m_{a}\cdot c_{a}\right)}{m_{g}^{2}}\right)^{2}+
	    				\left(\sigma_{C}\cdot\cfrac{T_{a}-T_{e}}{m_{g}}\right)^{2}
	    				\\
	    		\end{align*}

				Como não foi medida a temperatura de fusão, $T_{f}$, usaremos $T_{f}=\left(0\pm3\right)\unit{^{\circ}C}$.
				Os outros valores usados são os obtidos experimentalmente da Tabela
				2, assim obtem-se:

				\[
					L_{f}\pm\sigma_{L_{f}}=\left(6,4\pm0,8\right)\cdot10^{2}\unitfrac{cal}{g}
				\]


				Como o valor do calor latente de fusão da água é de $79.53\unitfrac{cal}{g}$,
				esta parte do nosso experimento também não obteve êxito.

			\subsubsection{Análise crítica dos resultados}
				Já o valor do calor latente de fusão do gelo não estar de acordo com o esperado,
				pode ser explicado por uma combinação das seguintes observações:
					\begin{enumerate}
						\item Como o experimento não foi realizado num refrigerador, é possível
						que parte da massa medida de gelo tenha derretido no processo antes
						de isso ser contabilizado, ou seja, o que foi apresentado como massa
						de gelo, provavelmente contém um híbrido entre gelo e água. A solução
						para isso seria realizar o experimento a uma temperatura inferior
						a da fusão do gelo.
						\item Além disso a temperatura de fusão não foi medida - mas isso não justifica
						o resultado completamente fora da faixa obtido, pois o erro usado cobriria este deslize.
						\item Quanto aos materiais, as mesmas justificativas usadas para os metais
						também vale aqui.
						\item Aqui além de não ser considerado o erro no calor específico da água
						também foi desconsiderado o calor específico do gelo.
					\end{enumerate}



	\section{Conclusão}



\end{document}
